\documentclass[11pt]{article}

\usepackage[a4paper,margin=2.6cm]{geometry}
\usepackage{amsmath,amssymb,amsthm}

\newcommand{\R}{\mathbb{R}}
\newcommand{\E}{\mathbb{E}}
\newcommand{\dd}{\,\mathrm{d}}

\title{Correction of Exercises in LaTeX\\Subject code: Corr\_exo}
\author{TODO: author name(s)}
\date{May 2026}

\begin{document}
\maketitle

\section*{Exercise 3.1: Existence of Implicit Schemes}

Let $\Psi:\R_+\times\R\times\R\to\R$ be Lipschitz in all variables. Denote by $L$
a Lipschitz constant of $\Psi$ with respect to its third variable. For fixed
$(t,x,h)$, define
\[
    \mathcal T(y)=x+h\Psi(t,x,y).
\]
If $hL<1$, then $\mathcal T$ is a contraction because
\[
    |\mathcal T(y_1)-\mathcal T(y_2)|
    \le hL |y_1-y_2|.
\]
By the Banach fixed point theorem, the equation
\[
    y=x+h\Psi(t,x,y)
\]
has a unique solution. A numerical method is Picard iteration:
\[
    y^{(0)}=x,\qquad
    y^{(k+1)}=x+h\Psi(t,x,y^{(k)}).
\]
For $hL<1$, this sequence converges geometrically to the unique solution
$s(t,x,h)$.

Now suppose $y=s(t,x,h)$ and $y=x+h\phi(t,x)$. Then
\[
    \phi(t,x)=\frac{s(t,x,h)-x}{h}
             =\Psi(t,x,s(t,x,h)).
\]
This is well-defined because $s(t,x,h)$ is unique. To prove that $s$ is
Lipschitz, let $s=s(t,x,h)$ and $\tilde s=s(\tilde t,\tilde x,h)$. Then
\[
\begin{aligned}
|s-\tilde s|
&\le |x-\tilde x|
   + hL\bigl(|t-\tilde t|+|x-\tilde x|+|s-\tilde s|\bigr).
\end{aligned}
\]
Thus, when $hL<1$,
\[
    |s-\tilde s|
    \le
    \frac{(1+hL)|x-\tilde x|+hL|t-\tilde t|}{1-hL}.
\]
So $s$ is Lipschitz. Since $\Psi$ is Lipschitz and $\phi(t,x)=\Psi(t,x,s(t,x,h))$,
the function $\phi$ is also Lipschitz for sufficiently small $h$.

\section*{Exercise 3.2: Euler-Lagrange Multipliers}

Let
\[
    L(\lambda,x,y)=f(x,y)+\lambda g(x,y).
\]
Since $y_0=Y(x_0)$, we have
\[
    g(x_0,y_0)=0.
\]
Also $\partial_y g(x_0,y_0)\ne 0$, so the equation
\[
    \partial_y L(\lambda_0,x_0,y_0)
    =\partial_y f(x_0,y_0)+\lambda_0\partial_y g(x_0,y_0)=0
\]
has the unique solution
\[
    \lambda_0
    =
    -\frac{\partial_y f(x_0,y_0)}{\partial_y g(x_0,y_0)}.
\]
Therefore
\[
    \partial_y L(\lambda_0,x_0,y_0)=0,
    \qquad
    \partial_\lambda L(\lambda_0,x_0,y_0)=g(x_0,y_0)=0.
\]

Differentiate the constraint $g(x,Y(x))=0$ at $x_0$:
\[
    \partial_x g(x_0,y_0)+\partial_y g(x_0,y_0)Y'(x_0)=0.
\]
Also
\[
    F'(x_0)=\partial_x f(x_0,y_0)+\partial_y f(x_0,y_0)Y'(x_0).
\]
Using $\partial_y f(x_0,y_0)=-\lambda_0\partial_y g(x_0,y_0)$ gives
\[
\begin{aligned}
F'(x_0)
&=\partial_x f(x_0,y_0)
  -\lambda_0\partial_y g(x_0,y_0)Y'(x_0)\\
&=\partial_x f(x_0,y_0)+\lambda_0\partial_x g(x_0,y_0)\\
&=\partial_x L(\lambda_0,x_0,y_0).
\end{aligned}
\]

For $f(x,y)=x+y$ and $g(x,y)=x^2+y^2-1$, the assumption
$\partial_y g(x_0,y_0)\ne0$ means $y_0\ne0$. We obtain
\[
    \lambda_0=-\frac{1}{2y_0}.
\]
Thus
\[
    F'(x_0)
    =\partial_x L(\lambda_0,x_0,y_0)
    =1+\lambda_0(2x_0)
    =1-\frac{x_0}{y_0}.
\]
This computes the derivative without explicitly solving for the local branch
$Y(x)$.

For the vector case, let
\[
    f:\R^n\times\R^n\to\R,\qquad
    g:\R^n\times\R^n\to\R^n,
\]
and define
\[
    L(\lambda,x,y)=f(x,y)+\lambda^\top g(x,y),
    \qquad \lambda\in\R^n.
\]
Assume $D_y g(x_0,y_0)$ is invertible. The multiplier is determined by
\[
    \nabla_y f(x_0,y_0)+D_y g(x_0,y_0)^\top\lambda_0=0,
\]
so
\[
    \lambda_0=-D_y g(x_0,y_0)^{-\top}\nabla_y f(x_0,y_0).
\]
Differentiating $g(x,Y(x))=0$ gives
\[
    D_x g(x_0,y_0)+D_y g(x_0,y_0)DY(x_0)=0.
\]
As in the scalar case, the terms containing $DY(x_0)$ cancel and
\[
    \nabla_x F(x_0)
    =
    \nabla_x f(x_0,y_0)+D_x g(x_0,y_0)^\top\lambda_0
    =
    \nabla_x L(\lambda_0,x_0,y_0).
\]

\section*{Exercise 3.3: Multi-Step SDE Scheme}

Consider
\[
    Y_{n+1}
    =
    Y_n+\frac32 a_nh-\frac12 a_{n-1}h+b_n\sqrt h\,\xi_n,
\]
where $\xi_n$ is independent of the filtration $\mathcal A_{t_n}$, has mean
$m$, and has variance $\sigma^2$.

For weak consistency, the conditional mean of the increment must match the drift
to first order. Since
\[
    \E[Y_{n+1}-Y_n\mid\mathcal A_{t_n}]
    =
    \frac32 a_nh-\frac12 a_{n-1}h+b_n\sqrt h\,m,
\]
the term $b_n\sqrt h\,m$ would be too large compared with the drift unless
\[
    m=0.
\]
With $m=0$, the Adams--Bashforth drift satisfies
\[
    \frac32 a_n-\frac12 a_{n-1}=a_n+o(1),
\]
because $a$ is smooth and $Y_{n-1}$ is close to $Y_n$ as $h\to0$.

For the second moment,
\[
    \E\!\left[(Y_{n+1}-Y_n)^2\mid\mathcal A_{t_n}\right]
    =
    b_n^2 h\,\E[\xi_n^2]+o(h).
\]
Since $m=0$, $\E[\xi_n^2]=\sigma^2$. Matching the diffusion coefficient requires
\[
    \sigma^2=1.
\]
Therefore the scheme is weakly consistent when
\[
    \boxed{m=0,\qquad \sigma^2=1.}
\]

The scheme is not strongly consistent in general with merely discrete variables
$\xi_n$. Strong consistency is pathwise and requires the noise increment to match
the Brownian increment in mean square. A standardized discrete variable may have
the correct first two moments, but it is not the same object as
$\Delta W_n/\sqrt h$. The mean-square difference between
$b_n\sqrt h\,\xi_n$ and $b_n\Delta W_n$ does not vanish at the local order needed
for strong consistency unless $\xi_n$ is chosen to be the actual normalized
Brownian increment. Thus the given weak scheme is generally not strongly
consistent.

\end{document}
